O desenvolvimento de veículos aéreos não tripulados (VANTs), em especial no contexto de projetos de aerodesign, exige o monitoramento rigoroso de variáveis cinemáticas e atmosféricas em tempo real para a validação empírica de modelos aerodinâmicos e estruturais. Contudo, as soluções comerciais de instrumentação em voo apresentam alto custo e arquitetura proprietária, limitando a customização e o acesso estudantil. Diante desta problemática, este trabalho propõe a concepção e implementação de um sistema de telemetria distribuído, de baixo custo e arquitetura aberta. A solução é composta por três nós integrados: um hardware embarcado baseado no microcontrolador ESP32 de duplo núcleo, responsável pela aquisição de dados inerciais, barométricos e de posicionamento global (GNSS), e transmissão via protocolo LoRa (915 MHz); um componente middleware terrestre que atua como roteador de protocolos, convertendo os datagramas de radiofrequência para pacotes UDP sobre uma rede Wi-Fi local; e uma estação de controle em solo estruturada sob o padrão Model-View-Controller (MVC), desenvolvida para a visualização gráfica contínua e o registro redundante dos logs de voo. O firmware embarcado adotou a Arquitetura Hexagonal, viabilizando ensaios de validação Hardware-in-the-Loop (HIL). Os resultados experimentais atestaram a estabilidade elétrica do circuito de aquisição sob carga e comprovaram a resiliência do enlace de comunicação, registrando apenas 1\% de perda de pacotes a 500 metros de distância sob condições não ideais de propagação. Adicionalmente, a estação de solo obteve êxito ao renderizar gráficos em uma taxa de atualização de 20 Hz sem bloqueios de processamento. Conclui-se que a arquitetura desenvolvida entrega uma plataforma telemétrica robusta, modular e de baixa latência, adequada para amparar ensaios de voo com alta confiabilidade.


\palavraschave{telemetria; vant; lora, esp32; sistema embarcado}
